\documentclass[12pt,a4paper]{article}
\usepackage{iftex}
\ifPDFTeX
  \usepackage[T2A]{fontenc}
  \usepackage[utf8]{inputenc}
\else
  \usepackage{fontspec}
  \setmainfont{Liberation Serif}
  \setmonofont{Liberation Mono}
\fi
\usepackage[russian]{babel}
\usepackage[margin=2cm]{geometry}
\usepackage{amsmath}
\usepackage{listings}

\pagestyle{empty}
\setlength{\parindent}{0pt}
\setlength{\parskip}{6pt}

\begin{document}

\begin{center}
{\Large \textbf{Контрольная работа N.1}}\\[4pt]
Вариант 0
\end{center}

\textbf{Время:} задания 1--3 по 10 минут; задача 4: 30 минут на вывод уравнения и 30 минут на решение в SymPy.\\
\textbf{Баллы:} задания 1--3 по 1 баллу, задача 4 -- 2 балла.

\hrule

\subsection*{Задание 1 (1 балл)}
Решить задачу Коши:
\[
(1+e^x)\,y\,y' = e^x, \qquad y(0)=1.
\]

\subsection*{Задание 2 (1 балл)}
Решить уравнение методом вариации постоянной и проверить ответ подстановкой в исходное уравнение:
\[
y' - \frac{2y}{x} = x^3 e^x, \qquad x>0.
\]

\subsection*{Задание 3 (1 балл)}
Решить уравнение Бернулли с помощью замены $y=u(x)\,v(x)$:
\[
y' + y = x\,y^2.
\]

\subsection*{Задание 4 (2 балла)}
Мяч массой $1$ кг бросили вертикально вверх со скоростью $20$ м/с с крыши здания высотой $50$ м.
На мяч действуют сила тяжести и сила сопротивления воздуха, пропорциональная скорости, с коэффициентом $0{,}5$ кг/с.
Ускорение свободного падения $g=10$ м/с$^2$.

\begin{enumerate}
\item[а)] На какой максимальной высоте над землёй окажется мяч?
\item[б)] Через какое время после броска и с какой скоростью мяч упадёт на землю?
\item[в)] К какой скорости стремилась бы скорость падающего мяча, если бы здание было сколь угодно высоким?
\item[г)] Постройте графики скорости $v(t)$ и высоты $x(t)$ от броска до падения.
\end{enumerate}

\textbf{Часть 1. Вывод дифференциального уравнения.}

Направим ось $x$ вверх, начало отсчёта --- на земле. Тогда начальные условия
\[
x(0)=50, \qquad x'(0)=20.
\]
На мяч действуют сила тяжести $-mg$ и сила сопротивления $-kx'$. Она всегда направлена против скорости, поэтому одна и та же формула описывает и подъём, и падение. Второй закон Ньютона:
\[
m x'' = -mg - k x'.
\]
При $m=1$, $k=0{,}5$, $g=10$ получаем
\[
x'' + 0{,}5\,x' + 10 = 0.
\]

\textbf{Часть 2. Решение.}
Решите задачу Коши для этого уравнения в SymPy с помощью \texttt{dsolve} и ответьте на вопросы а)--г).
Код решения лежит в отдельном файле \texttt{solution.py}.

Обратите внимание, что для решения необходимо использовать как \texttt{dsolve}, так и \texttt{nsolve} -- численный решатель.

\newpage
\section*{Ответы для самопроверки}

\textbf{1.} $y=\sqrt{1+2\ln\dfrac{1+e^x}{2}}$.

\textbf{2.} $y=x^2\big((x-1)e^x + C\big)$.

\textbf{3.} $y=\dfrac{1}{x+1-Ce^x}$, а также $y=0$.

\textbf{4.} $v(t)=40e^{-t/2}-20$, \quad $x(t)=130-80e^{-t/2}-20t$.\\
а) $x_{\max}=90-40\ln 2\approx 62{,}27$ м (при $t=2\ln 2\approx 1{,}39$ с);\\
б) $t\approx 6{,}33$ с, $v\approx -18{,}3$ м/с (вниз);\\
в) $-20$ м/с.

\end{document}
